\chapter{Affine Connections; Riemannian Connections}
\label{chap:dc-affine-connections-riemannian-connections}

Following the conventions of \cite{doCarmo}, $\mathcal{X}(M)$ denotes the space
of smooth vector fields on $M$, and $\mathcal{D}(M)$ denotes the ring of smooth
real-valued functions on $M$.

\section{Affine connections}

\begin{definition}[affine connection]
  \label{def:dc-ch2-2-1}
  \lean{Riemannian.IsAffineConnectionMap, Riemannian.AffineConnection.ofIsAffineConnectionMap, Riemannian.AffineConnection.isAffineConnectionMap}
  \leanok
  \dcref{ch2:2.1}
  An \emph{affine connection} on a differentiable manifold $M$ is a map
  \[
    \nabla\colon \mathcal{X}(M)\times\mathcal{X}(M)\longrightarrow\mathcal{X}(M),
    \qquad (X,Y)\longmapsto\nabla_XY,
  \]
  such that, for $X,Y,Z\in\mathcal{X}(M)$ and $f,g\in\mathcal{D}(M)$,
  \begin{enumerate}
    \item[(i)] $\nabla_{fX+gY}Z=f\nabla_XZ+g\nabla_YZ$;
    \item[(ii)] $\nabla_X(Y+Z)=\nabla_XY+\nabla_XZ$;
    \item[(iii)] $\nabla_X(fY)=f\nabla_XY+X(f)Y$.
  \end{enumerate}
\end{definition}

\begin{proposition}[covariant derivative along a curve]
  \label{prop:dc-ch2-2-2}
  \lean{Riemannian.AffineConnection.covDerivAlong,
    Riemannian.AffineConnection.covDerivAlong_eq_chart,
    Riemannian.AffineConnection.covDerivAlong_add,
    Riemannian.AffineConnection.covDerivAlong_smul,
    Riemannian.AffineConnection.covDerivAlong_comp_smoothVectorField,
    Riemannian.AffineConnection.covDerivAlong_unique,
    Riemannian.AffineConnection.covDerivAlong_unique_global,
    Riemannian.RiemannianMetric.leviCivita_covDerivAlong_eq}
  \uses{def:dc-ch2-2-1, lem:dc-ch2-2-2-coord, lem:dc-ch2-2-2-c-coord, lem:dc-ch2-2-2-c-bridge}
  \leanok
  \dcref{ch2:2.2}
  Let $M$ carry an affine connection $\nabla$, and let $c\colon I\to M$ be a
  differentiable curve. There is a unique operation assigning to every vector
  field $V$ along $c$ a vector field $DV/dt$ along $c$ such that
  \begin{enumerate}
    \item[(a)] $\frac{D}{dt}(V+W)=\frac{DV}{dt}+\frac{DW}{dt}$;
    \item[(b)] $\frac{D}{dt}(fV)=\frac{df}{dt}V+f\frac{DV}{dt}$ for every
      differentiable $f\colon I\to\mathbb{R}$;
    \item[(c)] if $V(t)=Y(c(t))$ for $Y\in\mathcal{X}(M)$, then
      $\frac{DV}{dt}=\nabla_{dc/dt}Y$.
  \end{enumerate}
  This operation is the \emph{covariant derivative along $c$}.
\end{proposition}
\begin{proof}
  \leanok
  \sketch
  In coordinates, write $V=\sum_jv^jX_j$, where
  $X_j=\frac{\partial}{\partial x_j}(c(t))$. Properties (a)--(c) and
  \cref{def:dc-ch2-2-1} force
  \[
    \frac{DV}{dt}
    =\sum_j\frac{dv^j}{dt}X_j
      +\sum_{i,j}\frac{dx_i}{dt}v^j\nabla_{X_i}X_j.\tag{1}
  \]
  Thus the operation is unique. Conversely, formula (1) defines it locally and
  gives (a)--(c); uniqueness makes the local definitions agree on chart
  overlaps.
\end{proof}

\begin{lemma}[covariant derivative in coordinates; properties (a), (b)]
  \label{lem:dc-ch2-2-2-coord}
  \lean{Riemannian.covariantDerivCoord,
    Riemannian.covariantDerivCoord_add,
    Riemannian.covariantDerivCoord_smul,
    Riemannian.covariantDerivCoord_unique}
  \uses{def:dc-ch2-3-7-christoffel}
  \leanok
  \dcref{ch2:2.2}
  Fix a Riemannian metric and a chart with Christoffel symbols $\Gamma^k_{ij}$.
  Set
  \[
    \Gamma(v,w)(y)=\sum_k\left(\sum_{i,j}\Gamma^k_{ij}(y)v^iw^j\right)e_k.
  \]
  For a coordinate curve $u\colon I\to\mathbb{R}^n$ and a coordinate vector
  field $V$ along it, formula (1) becomes
  \[
    \frac{DV}{dt}=\dot V+\Gamma(\dot u,V)(u).
  \]
  This operator is additive and satisfies
  $\frac{D}{dt}(fV)=\dot f\,V+f\frac{DV}{dt}$.
\end{lemma}

\begin{lemma}[property (c) in coordinates: induced fields]
  \label{lem:dc-ch2-2-2-c-coord}
  \lean{Riemannian.chartCovariantDeriv, Riemannian.covariantDerivCoord_induced}
  \uses{lem:dc-ch2-2-2-coord, def:dc-ch2-3-7-christoffel}
  \leanok
  \dcref{ch2:2.2}
  For a coordinate vector field $Y$ and a tangent vector $v$ at $y$, define
  \[
    (\nabla_vY)(y)=DY(y)v+\Gamma(v,Y(y))(y).
  \]
  If $V=Y\circ u$ along a coordinate curve $u$, then
  \[
    \frac{D}{dt}(Y\circ u)=(\nabla_{\dot u}Y)(u).
  \]
\end{lemma}
\begin{proof}
  \leanok
  \sketch
  The chain rule gives $(Y\circ u)'=DY(u)\dot u$, and the Christoffel terms in
  the two expressions coincide.
\end{proof}

\begin{lemma}[formula (10): the Koszul reduction on a coordinate frame]
  \label{lem:dc-ch2-2-2-c-formula10}
  \lean{Riemannian.christoffel_bridge_inner, Riemannian.christoffel_bridge_vector}
  \uses{lem:dc-ch2-3-6-koszul, def:dc-ch2-3-7-christoffel, thm:dc-ch2-3-6}
  \leanok
  \dcref{ch2:2.2}
  Let $\nabla$ be Levi-Civita and let $X_1,\dots,X_n$ be a coordinate frame at
  $p$. Suppose $[X_a,X_b](p)=0$ and
  $X_r\langle X_a,X_b\rangle(p)=\partial_rG_{ab}$, where
  $G_{ab}=\langle X_a,X_b\rangle$. Then
  \[
    \langle X_k,\nabla_{X_i}X_j\rangle
      =\sum_mG_{km}\Gamma^m_{ij},
    \qquad
    \nabla_{X_i}X_j=\sum_m\Gamma^m_{ij}X_m.\tag{10}
  \]
  \end{lemma}
\begin{proof}
    \leanok
    \sketch
    Apply the Koszul formula (9) from \cref{lem:dc-ch2-3-6-koszul} with
    $Y=X_i$, $X=X_j$, and $Z=X_k$. The bracket terms vanish, while the derivative
    terms give
    $\partial_iG_{kj}+\partial_jG_{ki}-\partial_kG_{ij}
      =2\sum_mG_{km}\Gamma^m_{ij}$.
    This proves the inner-product identity. Since the $X_k(p)$ form a basis and
    the metric is nondegenerate, it also proves the vector identity.
  \end{proof}

\begin{lemma}[property (c): chart-Christoffel bridge to the abstract connection]
  \label{lem:dc-ch2-2-2-c-bridge}
  \lean{Riemannian.exists_chartFrame_leviCivita_christoffel}
  \uses{lem:dc-ch2-2-2-c-coord, lem:dc-ch2-2-2-c-formula10, thm:dc-ch2-3-6, def:dc-ch2-3-7-christoffel}
  \leanok
  \dcref{ch2:2.2}
  The coordinate derivative in \cref{lem:dc-ch2-2-2-c-coord}, formed using the
  metric Christoffel symbols, is the coordinate expression of the abstract
  Levi-Civita connection from \cref{thm:dc-ch2-3-6}. At every point of a chart
  there are global smooth fields $X_1,\dots,X_n$ agreeing locally with its
  coordinate frame and satisfying
  \[
    \nabla_{X_i}X_j=\sum_m\Gamma^m_{ij}X_m.
  \]
  \end{lemma}
\begin{proof}
    \leanok
    \sketch
    Extend the coordinate frame fields globally without changing their germs at
    the chosen point. Locally they are pullbacks of constant coordinate fields,
    so their brackets vanish, and the directional derivatives of their Gram
    entries are the corresponding partial derivatives. Formula (10) from
    \cref{lem:dc-ch2-2-2-c-formula10} then identifies the coordinate and abstract
    Christoffel symbols, establishing property (c) of
    \cref{prop:dc-ch2-2-2}.
  \end{proof}

\begin{remark}[Remark 2.3]
  \label{rem:dc-ch2-2-3}
  \lean{Riemannian.AffineConnection.cov_congr_apply_left, Riemannian.AffineConnection.cov_congr_apply_eventuallyEq_right, Riemannian.AffineConnection.cov_smulSum_apply}
  \uses{rem:dc-ch0-5-7, def:dc-ch2-2-1}
  \leanok
  \dcref{ch2:2.3}
  An affine connection is local. In coordinates, if
  $X=\sum_ix_iX_i$, $Y=\sum_jy_jX_j$, and
  $\nabla_{X_i}X_j=\sum_k\Gamma_{ij}^kX_k$, then
  \[
    \nabla_XY
      =\sum_{i,j}x_iy_j\nabla_{X_i}X_j+\sum_jX(y_j)X_j
      =\sum_k\left(\sum_{i,j}x_iy_j\Gamma_{ij}^k+X(y_k)\right)X_k.
  \]
  Hence $\nabla_XY(p)$ depends only on $X(p)$ and the germ of $Y$ at $p$.
\end{remark}

\begin{remark}[Remark 2.4]
  \label{rem:dc-ch2-2-4}
  \lean{Riemannian.AffineConnection.covDerivAlong,
    Riemannian.Geodesic.continuousAt_and_hasGeodesicEquationAt_iff_leviCivita_covDerivAlong_velocity_eq_zero,
    Riemannian.covariantDerivCoord,
    Riemannian.covariantDerivCoord_add,
    Riemannian.covariantDerivCoord_smul}
  \uses{prop:dc-ch2-2-2}
  \leanok
  \dcref{ch2:2.4}
  An affine connection supplies a derivative of vector fields along curves, and
  therefore defines the acceleration $D(dc/dt)/dt$ of a curve in $M$.
\end{remark}

\begin{definition}[parallel vector field]
  \label{def:dc-ch2-2-5}
  \lean{Riemannian.IsParallelAlongWithinOn,
    Riemannian.IsIntrinsicPiecewiseParallelAlongOn,
    Riemannian.IsParallelCoord,
    Riemannian.IsParallelAlong,
    Riemannian.IsParallelAlongOn,
    Riemannian.isParallelAlongWithinOn_iff_isParallelFieldAlongOn,
    Riemannian.RiemannianMetric.leviCivita_isParallelFieldAlongOn_iff_isParallelAlong}
  \uses{lem:dc-ch2-2-2-coord}
  \leanok
  \dcref{ch2:2.5}
  A vector field $V$ along $c\colon I\to M$ is \emph{parallel} if
  $DV/dt=0$ on $I$. In coordinates this is the linear system
  \[
    \dot V=-\Gamma(\dot u,V)(u).
  \]
\end{definition}

\begin{lemma}[the parallel-transport coefficient is a bounded linear operator]
  \label{lem:dc-ch2-2-6-coeff}
  \lean{Riemannian.chartChristoffelContractionRight, Riemannian.lipschitzWith_neg_chartChristoffelContraction}
  \uses{def:dc-ch2-2-5}
  \leanok
  \dcref{ch2:2.6}
  For a chart velocity $v$ and base point $y$, the map
  $w\mapsto\Gamma(v,w)(y)$ is a bounded linear operator $A(v,y)$ on the
  finite-dimensional model space. Thus
  $w\mapsto-\Gamma(\dot u(t),w)(u(t))$ is globally Lipschitz with constant
  $\lVert A(\dot u(t),u(t))\rVert$.
\end{lemma}

\begin{proposition}[parallel transport]
  \label{prop:dc-ch2-2-6}
  \lean{Riemannian.exists_unique_isParallelAlongWithinOn,
    Riemannian.exists_unique_isIntrinsicPiecewiseParallelAlongOn,
    Riemannian.intrinsicPiecewiseParallelTransportTangentEquiv,
    Riemannian.intrinsicPiecewiseParallelTransportTangentEquiv_eq_terminal,
    Riemannian.metricInner_intrinsicPiecewiseParallelTransportTangentEquiv}
  \uses{def:dc-ch2-2-5, lem:dc-ch2-2-6-exist, lem:dc-ch2-2-6-unique, lem:dc-ch2-2-6-transport}
  \leanok
  \dcref{ch2:2.6}
  Let $c\colon I\to M$ be differentiable, $t_0\in I$, and
  $V_0\in T_{c(t_0)}M$. There is a unique parallel vector field $V$ along $c$
  with $V(t_0)=V_0$. Its values are the \emph{parallel transports} of $V_0$
  along $c$.
\end{proposition}
\begin{proof}
  \uses{lem:dc-ch2-2-6-exist, lem:dc-ch2-2-6-unique, lem:dc-ch2-2-6-transport}
  \sketch
  In one chart, writing $V=\sum_jv^jX_j$ reduces parallelism to
  \[
    0=\frac{dv^k}{dt}
      +\sum_{i,j}\Gamma_{ij}^kv^j\frac{dx_i}{dt},
      \qquad k=1,\dots,n.\tag{2}
  \]
  This linear system has a unique solution for each initial value. On a compact
  subinterval, cover the curve by finitely many coordinate neighborhoods and
  glue the local solutions using uniqueness. These solutions extend throughout
  $I$.
\end{proof}

\begin{lemma}[existence of parallel transport]
  \label{lem:dc-ch2-2-6-exist}
  \lean{Riemannian.exists_isParallelCoord_Icc,
    Riemannian.exists_isParallelAlongWithinOn,
    Riemannian.exists_isParallelAlongPresentation}
  \uses{def:dc-ch2-2-5, lem:dc-ch2-2-6-coeff}
  \leanok
  \dcref{ch2:2.6}
  On a compact interval $[a,b]$, every initial coordinate vector $V_0$ extends
  to a solution of
  \[
    \dot V=-\Gamma(\dot u,V)(u),\qquad V(a)=V_0.
  \]
  It suffices that $u$ is $C^1$ and the Christoffel symbols are continuous.
\end{lemma}
\begin{proof}
  \sketch
  By \cref{lem:dc-ch2-2-6-coeff}, this is a linear ODE
  $\dot V=A(t)V$ with continuous bounded coefficients. Local existence and a
  finite subdivision of $[a,b]$ give a global solution.
\end{proof}

\begin{lemma}[uniqueness of parallel transport]
  \label{lem:dc-ch2-2-6-unique}
  \lean{Riemannian.isParallelCoord_eqOn_Icc,
    Riemannian.isParallelSol_eqOn_Icc,
    Riemannian.IsParallelAlongWithinOn.eqOn_of_eq_at,
    Riemannian.eqOn_of_two_intrinsicPresentations,
    Riemannian.IsIntrinsicPiecewiseParallelAlongOn.eqOn}
  \uses{def:dc-ch2-2-5, lem:dc-ch2-2-2-coord, lem:dc-ch2-2-6-coeff}
  \leanok
  \dcref{ch2:2.6}
  Two parallel coordinate vector fields along the same coordinate curve on
  $[a,b]$ that agree at $a$ agree throughout $[a,b]$.
\end{lemma}
\begin{proof}
  \sketch
  Their difference solves the homogeneous linear system (2) with zero initial
  value. The Lipschitz estimate from \cref{lem:dc-ch2-2-6-coeff} and
  Grönwall's inequality imply that it vanishes.
\end{proof}

\begin{lemma}[parallel transport is a linear isomorphism]
  \label{lem:dc-ch2-2-6-transport}
  \lean{Riemannian.parallelTransport,
    Riemannian.exists_parallelTransport_spec,
    Riemannian.intrinsicPiecewiseParallelTransportTangentEquiv,
    Riemannian.intrinsicPiecewiseParallelTransportTangentEquiv_eq_terminal,
    Riemannian.metricInner_intrinsicPiecewiseParallelTransportTangentEquiv}
  \uses{lem:dc-ch2-2-6-exist, lem:dc-ch2-2-6-unique}
  \leanok
  \dcref{ch2:2.6}
  On $[a,b]$, the endpoint map
  \[
    P_c\colon T_{c(a)}M\longrightarrow T_{c(b)}M,
    \qquad V_0\longmapsto V(b),
  \]
  where $V$ is the parallel field with $V(a)=V_0$, is a linear isomorphism.
  Its inverse is parallel transport along the reversed curve.
\end{lemma}
\begin{proof}
  \sketch
  Existence and uniqueness make the map well-defined. Superposition for the
  linear system (2) gives linearity, and backward uniqueness gives injectivity
  and the stated inverse.
\end{proof}

\section{Riemannian connections}

\begin{definition}[compatible with the metric]
  \label{def:dc-ch2-3-1}
  \lean{Riemannian.CompatibleGen,
    Riemannian.AffineConnection.PreservesMetricUnderParallelTransport,
    Riemannian.AffineConnection.PreservesMetricAlongOn}
  \leanok
  \dcref{ch2:3.1}
  An affine connection $\nabla$ on a Riemannian manifold is \emph{compatible
  with the metric} if the inner product of any two parallel vector fields along
  any smooth curve is constant.
\end{definition}

\begin{proposition}[product rule characterization]
  \label{prop:dc-ch2-3-2}
  \lean{Riemannian.compatibleGen_iff_productRuleGen,
    Riemannian.AffineConnection.HasMetricProductRuleAlong,
    Riemannian.AffineConnection.preservesMetricUnderParallelTransport_iff_hasMetricProductRuleAlong,
    Riemannian.hasDerivAt_metricInner_along_of_mdifferentiableAt,
    Riemannian.RiemannianMetric.leviCivitaConnection_hasMetricProductRuleAlong,
    Riemannian.metricInner_const_of_isParallelAlong,
    Riemannian.metricInner_eq_endpoints_of_isParallelAlongWithinOn}
  \uses{def:dc-ch2-3-1, prop:dc-ch2-2-6, lem:dc-ch2-3-2-coord}
  \leanok
  \dcref{ch2:3.2}
  A connection $\nabla$ is compatible with the metric if and only if, for all
  vector fields $V,W$ along every differentiable curve $c\colon I\to M$,
  \[
    \frac{d}{dt}\langle V,W\rangle
      =\left\langle\frac{DV}{dt},W\right\rangle
       +\left\langle V,\frac{DW}{dt}\right\rangle,
    \qquad t\in I.\tag{3}
  \]
\end{proposition}
\begin{proof}
  \leanok
  \sketch
  Formula (3) immediately makes inner products of parallel fields constant.
  Conversely, parallel-transport an orthonormal basis at one point along $c$;
  compatibility keeps it orthonormal. Expanding
  $V=\sum_iv^iP_i$ and $W=\sum_iw^iP_i$ in this parallel frame reduces (3) to
  the ordinary product rule for $\sum_iv^iw^i$.
\end{proof}

\begin{corollary}[compatibility on vector fields]
  \label{cor:dc-ch2-3-3}
  \lean{Riemannian.AffineConnection.IsMetricCompatible,
    Riemannian.AffineConnection.isMetricCompatible_iff_hasMetricProductRuleAlong,
    Riemannian.AffineConnection.preservesMetricUnderParallelTransport_iff_isMetricCompatible}
  \uses{prop:dc-ch2-3-2}
  \leanok
  \dcref{ch2:3.3}
  A connection $\nabla$ is compatible with the metric if and only if
  \[
    X\langle Y,Z\rangle
      =\langle\nabla_XY,Z\rangle+\langle Y,\nabla_XZ\rangle,
    \qquad X,Y,Z\in\mathcal{X}(M).\tag{4}
  \]
\end{corollary}
\begin{proof}
  \sketch
  Evaluate (3) at a point along a curve whose velocity there is $X$. This gives
  (4). Conversely, apply (4) pointwise to the velocity field of a curve to
  recover (3).
\end{proof}

\begin{lemma}[metric compatibility of $D/dt$, coordinate form]
  \label{lem:dc-ch2-3-2-coord}
  \lean{Riemannian.CompatibleGen,
    Riemannian.compatibleGen_iff_productRuleGen,
    Riemannian.hasDerivAt_chartMetricInner_along}
  \uses{lem:dc-ch2-2-2-coord, def:dc-ch2-3-7-christoffel}
  \leanok
  \dcref{ch2:3.2}
  In a chart, write
  $\langle a,b\rangle_y=\sum_{i,j}g_{ij}(y)a^ib^j$ and
  $DV/dt=\dot V+\Gamma(\dot u,V)(u)$. Then coordinate vector fields $V,W$
  along $u$ satisfy
  \[
    \frac{d}{dt}\langle V,W\rangle
      =\left\langle\frac{DV}{dt},W\right\rangle
       +\left\langle V,\frac{DW}{dt}\right\rangle.
  \]
\end{lemma}
\begin{proof}
  \leanok
  \uses{def:dc-ch2-3-7-christoffel}
  \sketch
  Differentiate $\sum_{i,j}g_{ij}(u)V^iW^j$ and use
  \[
    \partial_kg_{ij}
      =\sum_m\left(g_{mj}\Gamma^m_{ki}+g_{im}\Gamma^m_{kj}\right),
  \]
  which follows from \cref{def:dc-ch2-3-7-christoffel}. The two Christoffel
  sums combine respectively with $\dot V$ and $\dot W$ to give the two
  covariant derivatives.
\end{proof}

\begin{lemma}[parallel fields have constant inner product, coordinate form]
  \label{lem:dc-ch2-3-1-coord}
  \lean{Riemannian.CompatibleGen,
    Riemannian.hasDerivAt_chartMetricInner_eq_zero_of_isParallelCoord}
  \uses{def:dc-ch2-2-5, def:dc-ch2-3-1, lem:dc-ch2-3-2-coord}
  \leanok
  \dcref{ch2:3.1}
  If $V$ and $W$ are parallel coordinate vector fields along a coordinate
  curve, then $\langle V,W\rangle$ is constant.
\end{lemma}
\begin{proof}
  \leanok
  \uses{lem:dc-ch2-3-2-coord}
  \sketch
  In \cref{lem:dc-ch2-3-2-coord}, both covariant derivatives vanish, so
  $\frac{d}{dt}\langle V,W\rangle=0$.
\end{proof}

\begin{definition}[symmetric connection]
  \label{def:dc-ch2-3-4}
  \lean{Riemannian.AffineConnection.IsSymmetric}
  \uses{def:dc-ch2-2-1}
  \leanok
  \dcref{ch2:3.4}
  An affine connection $\nabla$ is \emph{symmetric} if
  \[
    \nabla_XY-\nabla_YX=[X,Y]
    \quad\text{for all }X,Y\in\mathcal{X}(M).\tag{5}
  \]
\end{definition}

\begin{remark}[Remark 3.5]
  \label{rem:dc-ch2-3-5}
  \lean{Riemannian.chartChristoffel_symm}
  \leanok
  \dcref{ch2:3.5}
  In a coordinate frame $X_i=\partial/\partial x_i$, symmetry is equivalent to
  \[
    \nabla_{X_i}X_j-\nabla_{X_j}X_i=[X_i,X_j]=0,
    \qquad
    \Gamma_{ij}^k=\Gamma_{ji}^k.\tag{5'}
  \]
\end{remark}

\begin{lemma}[Koszul formula]
  \label{lem:dc-ch2-3-6-koszul}
  \lean{Riemannian.AffineConnection.koszul_formula}
  \uses{def:dc-ch2-3-4, cor:dc-ch2-3-3}
  \leanok
  \dcref{ch2:3.6}
  If $\nabla$ is symmetric and metric-compatible, then, for all
  $X,Y,Z\in\mathcal{X}(M)$,
  \[
    \begin{aligned}
    2\langle Z,\nabla_YX\rangle
      ={}&X\langle Y,Z\rangle+Y\langle Z,X\rangle-Z\langle X,Y\rangle\\
       &-\langle[X,Z],Y\rangle-\langle[Y,Z],X\rangle
        -\langle[X,Y],Z\rangle.
    \end{aligned}\tag{9}
  \]
\end{lemma}
\begin{proof}
  \leanok
  \uses{def:dc-ch2-3-4, cor:dc-ch2-3-3}
  \sketch
  Apply (4) to the three cyclic orderings of $X,Y,Z$. Add the first two
  identities, subtract the third, and replace differences of covariant
  derivatives using symmetry (5). Symmetry of the metric gives (9).
\end{proof}

\begin{lemma}[tangent-vector extension]
  \label{lem:dc-ch2-tangent-extension}
  \lean{Riemannian.exists_smoothVectorField_eq}
  \leanok
  On a finite-dimensional, $\sigma$-compact, Hausdorff smooth manifold, every
  $v\in T_pM$ is the value $Z(p)$ of a global smooth vector field
  $Z\in\mathcal{X}(M)$.
\end{lemma}
\begin{proof}
  \leanok
  \sketch
  In a bundle trivialization near $p$, take the locally constant section with
  value $v$ at $p$; away from $p$, use the zero section. A partition of unity
  glues these local sections to the required global field.
\end{proof}

\begin{lemma}[uniqueness of the Levi-Civita connection]
  \label{lem:dc-ch2-3-6-unique}
  \lean{Riemannian.AffineConnection.leviCivita_unique'}
  \uses{lem:dc-ch2-3-6-koszul, lem:dc-ch2-tangent-extension}
  \leanok
  \dcref{ch2:3.6}
  On a finite-dimensional, $\sigma$-compact, Hausdorff Riemannian manifold,
  there is at most one affine connection that is symmetric and
  metric-compatible.
\end{lemma}
\begin{proof}
  \leanok
  \uses{lem:dc-ch2-3-6-koszul, lem:dc-ch2-tangent-extension}
  \sketch
  For two such connections, (9) gives equal inner products
  $\langle Z,(\nabla_1)_YX\rangle=\langle Z,(\nabla_2)_YX\rangle$ for every
  global field $Z$. By \cref{lem:dc-ch2-tangent-extension}, this tests against
  every tangent vector at each point; nondegeneracy therefore gives
  $\nabla_1=\nabla_2$.
\end{proof}

\begin{lemma}[well-definedness of the Koszul functional]
  \label{lem:dc-ch2-3-6-welldef}
  \lean{Riemannian.RiemannianMetric.koszulRHS, Riemannian.RiemannianMetric.koszulRHS_add_right, Riemannian.RiemannianMetric.koszulRHS_smul_right}
  \uses{lem:dc-ch2-3-6-koszul}
  \leanok
  \dcref{ch2:3.6}
  For fixed $X,Y\in\mathcal{X}(M)$, let
  \[
    F_{Y,X}(Z)
      =X\langle Y,Z\rangle+Y\langle Z,X\rangle-Z\langle X,Y\rangle
       -\langle[X,Z],Y\rangle-\langle[Y,Z],X\rangle
       -\langle[X,Y],Z\rangle.
  \]
  Then $F_{Y,X}$ is $\mathcal{D}(M)$-linear in $Z$. Hence
  $Z\mapsto\frac12F_{Y,X}(Z)$ is tensorial and has a metric Riesz dual, the
  candidate $\nabla_YX$.
\end{lemma}
\begin{proof}
  \leanok
  \uses{lem:dc-ch2-3-6-koszul}
  \sketch
  Additivity is termwise. For $f\in\mathcal{D}(M)$, expand
  $X(fh)=fX(h)+X(f)h$ and
  $[X,fZ]=f[X,Z]+X(f)Z$. The derivative-of-$f$ terms cancel by symmetry of the
  metric, leaving $F_{Y,X}(fZ)=fF_{Y,X}(Z)$.
\end{proof}

\begin{lemma}[the Koszul-dual connection is an affine connection]
  \label{lem:dc-ch2-3-6-affine}
  \lean{Riemannian.RiemannianMetric.koszulRHS_add_left, Riemannian.RiemannianMetric.koszulRHS_leibniz_left, Riemannian.RiemannianMetric.koszulRHS_add_middle, Riemannian.RiemannianMetric.koszulRHS_smul_middle}
  \uses{def:dc-ch2-2-1, lem:dc-ch2-3-6-welldef}
  \leanok
  \dcref{ch2:3.6}
  The candidate $\nabla_YX=\frac12(F_{Y,X})^\sharp$ satisfies
  \begin{align*}
    F_{Y,X_1+X_2}&=F_{Y,X_1}+F_{Y,X_2},\\
    F_{Y,fX}&=fF_{Y,X}+2Y(f)\langle X,Z\rangle,\\
    F_{Y_1+Y_2,X}&=F_{Y_1,X}+F_{Y_2,X},\\
    F_{fY,X}&=fF_{Y,X}.
  \end{align*}
  Therefore it is additive and $\mathcal{D}(M)$-linear in the direction field,
  and additive with the Leibniz rule in the differentiated field; hence it is an
  affine connection.
\end{lemma}
\begin{proof}
  \leanok
  \uses{lem:dc-ch2-3-6-koszul}
  \sketch
  Substitute the indicated sums and scalar multiples into the six terms of
  $F$. Additivity is termwise. The Leibniz rules for derivatives and brackets
  produce the displayed $2Y(f)\langle X,Z\rangle$ in the differentiated slot;
  in the direction slot all derivative-of-$f$ terms cancel in pairs. Together
  with \cref{lem:dc-ch2-3-6-welldef}, these identities give the connection
  axioms.
\end{proof}

\begin{lemma}[the Koszul-dual connection is Levi-Civita]
  \label{lem:dc-ch2-3-6-exist}
  \lean{Riemannian.AffineConnection.isLeviCivita_of_koszulDual, Riemannian.koszulDual_isSymmetric, Riemannian.koszulDual_isMetricCompatible}
  \uses{def:dc-ch2-2-1, def:dc-ch2-3-4, lem:dc-ch2-3-6-koszul, lem:dc-ch2-3-6-welldef, lem:dc-ch2-tangent-extension}
  \leanok
  \dcref{ch2:3.6}
  On a finite-dimensional, $\sigma$-compact, Hausdorff Riemannian manifold,
  suppose an affine connection satisfies
  \[
    2\langle\nabla_YX,Z\rangle=F_{Y,X}(Z)
  \]
  for every $X,Y,Z$. Then it is symmetric and metric-compatible.
\end{lemma}
\begin{proof}
  \leanok
  \uses{lem:dc-ch2-3-6-koszul, lem:dc-ch2-tangent-extension}
  \sketch
  The identities
  \[
    F_{Y,X}(Z)-F_{X,Y}(Z)=2\langle[X,Y],Z\rangle
  \]
  and
  \[
    F_{Y,X}(Z)+F_{Z,X}(Y)=2X\langle Y,Z\rangle
  \]
  follow by cancellation. The first, together with
  \cref{lem:dc-ch2-tangent-extension} and nondegeneracy, gives
  $\nabla_XY-\nabla_YX=[X,Y]$. The second gives the product rule (4), hence
  metric compatibility.
\end{proof}

\begin{theorem}[Levi-Civita]
  \label{thm:dc-ch2-3-6}
  \lean{Riemannian.RiemannianMetric.exists_unique_isLeviCivita}
  \uses{def:dc-ch2-3-4, cor:dc-ch2-3-3, lem:dc-ch2-3-6-koszul, lem:dc-ch2-3-6-unique, lem:dc-ch2-3-6-welldef, lem:dc-ch2-3-6-exist}
  \leanok
  \dcref{ch2:3.6}
  Every Riemannian manifold admits a unique affine connection $\nabla$ that is
  \begin{enumerate}
    \item[(a)] symmetric;
    \item[(b)] compatible with the Riemannian metric.
  \end{enumerate}
\end{theorem}
\begin{proof}
  \leanok
  \uses{lem:dc-ch2-3-6-unique, lem:dc-ch2-3-6-exist}
  \sketch
  Metric compatibility applied cyclically gives
  \[
    X\langle Y,Z\rangle
      =\langle\nabla_XY,Z\rangle+\langle Y,\nabla_XZ\rangle,\tag{6}
  \]
  \[
    Y\langle Z,X\rangle
      =\langle\nabla_YZ,X\rangle+\langle Z,\nabla_YX\rangle,\tag{7}
  \]
  \[
    Z\langle X,Y\rangle
      =\langle\nabla_ZX,Y\rangle+\langle X,\nabla_ZY\rangle.\tag{8}
  \]
  Adding (6) and (7), subtracting (8), and using symmetry gives
  \[
    \begin{aligned}
    \langle Z,\nabla_YX\rangle=\tfrac12\{&X\langle Y,Z\rangle
      +Y\langle Z,X\rangle-Z\langle X,Y\rangle\\
      &-\langle[X,Z],Y\rangle-\langle[Y,Z],X\rangle
       -\langle[X,Y],Z\rangle\}.\tag{9}
    \end{aligned}
  \]
  Thus uniqueness follows from \cref{lem:dc-ch2-3-6-unique}. For existence,
  take the metric Riesz dual of the right-hand side of (9); its
  well-definedness and Levi-Civita properties are
  \cref{lem:dc-ch2-3-6-welldef,lem:dc-ch2-3-6-exist}.
\end{proof}

\begin{definition}[Christoffel symbols]
  \label{def:dc-ch2-3-7-christoffel}
  \lean{Riemannian.chartChristoffel}
  \uses{def:dc-ch1-2-1}
  \leanok
  \dcref{ch2:3.7}
  In coordinates, let $g_{ij}=\langle X_i,X_j\rangle$ and let $(g^{km})$ be
  the inverse matrix. The \emph{Christoffel symbols of the second kind} are
  \[
    \Gamma_{ij}^m
      =\tfrac12\sum_k\left(
        \frac{\partial g_{jk}}{\partial x_i}
       +\frac{\partial g_{ki}}{\partial x_j}
       -\frac{\partial g_{ij}}{\partial x_k}\right)g^{km}.\tag{10}
  \]
\end{definition}

\begin{remark}[the Levi-Civita connection and its coefficients]
  \label{rem:dc-ch2-3-7}
  \lean{Riemannian.chartChristoffel, Riemannian.chartChristoffel_symm, Riemannian.chartGram_christoffel_contraction, Riemannian.christoffel_bridge_inner, Riemannian.christoffel_bridge_vector, Riemannian.exists_chartFrame_leviCivita_christoffel}
  \uses{thm:dc-ch2-3-6, def:dc-ch2-3-7-christoffel}
  \leanok
  \dcref{ch2:3.7}
  The connection of \cref{thm:dc-ch2-3-6} is the \emph{Levi-Civita} or
  \emph{Riemannian connection}. Its coordinate coefficients, defined by
  $\nabla_{X_i}X_j=\sum_k\Gamma_{ij}^kX_k$, satisfy
  \[
    \sum_\ell\Gamma_{ij}^\ell g_{\ell k}
      =\tfrac12\left(
        \frac{\partial g_{jk}}{\partial x_i}
       +\frac{\partial g_{ki}}{\partial x_j}
       -\frac{\partial g_{ij}}{\partial x_k}\right),
  \]
  and hence agree with \cref{def:dc-ch2-3-7-christoffel}. Moreover,
  \[
    \frac{DV}{dt}
      =\sum_k\left(\frac{dv^k}{dt}
        +\sum_{i,j}\Gamma_{ij}^kv^j\frac{dx_i}{dt}\right)X_k.
  \]
  In Euclidean coordinates, $\Gamma_{ij}^k=0$, so this is the ordinary
  derivative.
\end{remark}

\begin{proposition}[Riemannian parallel transport along a $C^1$ curve]
  \label{prop:dc-ch2-2-6-riemannian}
  \dcref{ch2:2.6}
  \uses{def:dc-ch2-2-5}
  \lean{Riemannian.Variation.exists_parallelCovariantFieldAlongOn_at}
  \leanok
  Let $M$ be a positive-dimensional boundaryless Riemannian manifold, let
  $\gamma:\mathbb{R}\to M$ be a $C^1$ curve, and let $a<b$. For every
  $t_0\in[a,b]$ and $V_0\in T_{\gamma(t_0)}M$, there exists a parallel vector
  field $V$ along $\gamma$ on $[a,b]$ such that $V(t_0)=V_0$. It is unique on
  $[a,b]$ among parallel fields having that value at $t_0$.
\end{proposition}
\begin{proof}
  \uses{lem:dc-ch2-2-6-exist, lem:dc-ch2-2-6-unique, lem:dc-ch2-2-6-transport}
  \leanok
  Starting at $a$, solve the parallelism system in a coordinate neighborhood.
  Compactness of $[a,b]$ gives finitely many such neighborhoods along the curve,
  and uniqueness makes the local solutions agree on every overlap; hence each
  $X\in T_{\gamma(a)}M$ determines a unique parallel field $V_X$ on $[a,b]$.
  Evaluation at $t_0$ defines a linear map
  $P_{a,t_0}:T_{\gamma(a)}M\to T_{\gamma(t_0)}M$. If
  $P_{a,t_0}(X)=P_{a,t_0}(Y)$, uniqueness applied from $t_0$ gives
  $V_X=V_Y$ on $[a,b]$, and therefore $X=Y$. Thus $P_{a,t_0}$ is injective;
  since its domain and codomain have the same finite dimension, it is an
  isomorphism. Set $X=P_{a,t_0}^{-1}(V_0)$ and $V=V_X$. Then
  $V(t_0)=V_0$, and a final application of uniqueness at $t_0$ proves the
  asserted uniqueness of $V$ on $[a,b]$.
\end{proof}
